\chapter{\MakeUppercase{Fundamentação e Trabalhos Relacionados}}\label{sec:fundamentacao}

Este capítulo apresenta os conceitos técnicos sobre os quais o AgroVisible se apoia, sintetiza o mapeamento sistemático da literatura realizado para posicionar o trabalho e caracteriza a lacuna científica que ele pretende preencher.

\section{Fundamentação Técnica}\label{sec2:fund-tecnica}

\subsection{Redes neurais recorrentes para séries temporais}

Redes neurais recorrentes processam sequências mantendo um estado interno que é atualizado a cada passo de tempo, o que as torna adequadas a séries temporais em que a ordem das observações importa. A arquitetura recorrente clássica, porém, tem dificuldade em reter informação por muitos passos, porque os gradientes que propagam o erro pela sequência se atenuam ou explodem com o comprimento da janela.

A rede de memória de curto e longo prazo (LSTM, \textit{Long Short-Term Memory}), proposta por \citeonline{Hochreiter1997}, resolve esse problema com uma célula de memória controlada por três portões (entrada, esquecimento e saída), que aprendem quanto da informação nova incorporar, quanto do estado anterior descartar e quanto do estado atual expor. Isso permite que a rede capture dependências de prazo mais longo, como a relação entre um acumulado de chuva de 90 dias e o retorno do preço seis meses depois.

Em previsão de preços de commodities, a LSTM é a arquitetura mais empregada na literatura recente, isoladamente ou combinada com camadas convolucionais, decomposição de sinal ou mecanismos de atenção \cite{K2023, Destriana2025, Huang2025}. Dois cuidados são recorrentes nesses trabalhos e adotados aqui: prever retornos, e não níveis, para evitar que a rede aprenda apenas a tendência; e avaliar contra referências ingênuas, porque séries de preço financeiro são próximas de um passeio aleatório e uma rede que apenas reproduz o último valor observado pode exibir métricas enganosamente boas.

\subsection{Processos de decisão de Markov e a equação de Bellman}

Um processo de decisão de Markov (MDP) é um modelo de decisão sequencial sob incerteza definido por um conjunto de estados, um conjunto de ações, uma função de transição que dá a probabilidade do próximo estado dado o estado e a ação atuais, e uma função de recompensa \cite{Puterman1994}. A propriedade de Markov exige que a transição dependa apenas do estado presente, e não de todo o histórico.

Uma política é uma regra que associa a cada estado uma ação. A política ótima maximiza a recompensa acumulada esperada, e sua função de valor satisfaz a equação de Bellman \cite{Bellman1957}: o valor de um estado é o máximo, sobre as ações, da recompensa imediata mais o valor esperado descontado do estado seguinte. Em horizonte finito, a equação é resolvida por indução retroativa, partindo do último período, em que o valor é conhecido, e recuando até o primeiro. Quando o espaço de estados é discreto e o horizonte finito, o problema equivale a encontrar o caminho de maior valor esperado em um grafo direcionado acíclico organizado em camadas temporais.

O problema de comercialização pós-colheita ajusta-se naturalmente a esse formalismo. O estado é o preço corrente, a ação é vender ou manter o estoque, a transição é a distribuição do preço futuro, e a recompensa é a receita líquida dos custos de carregamento. A solução por indução retroativa produz não apenas a ação recomendada hoje, mas a regra completa de decisão para cada preço possível em cada mês, o que permite comunicar ao produtor um gatilho de venda em vez de uma única recomendação.

\subsection{Custo de carregamento, custo de oportunidade e base}

Manter café estocado tem três custos. O primeiro é o custo físico de armazenagem, cobrado por armazéns gerais e cooperativas por saca e por mês, acrescido da quebra técnica, a perda de peso por secagem e manuseio. O segundo é o custo de oportunidade do capital: o produtor que estoca deixa de receber hoje um valor que poderia ser aplicado à taxa livre de risco, no Brasil referenciada pela Selic. Esse custo é frequentemente ignorado pelo produtor, mas em um ambiente de juros de dois dígitos ele supera a armazenagem física por larga margem. O terceiro é o custo de produção já incorrido, medido pelo COE e pelo COT da CONAB, que define a margem da safra mas não altera a decisão marginal entre vender e estocar, por já ser um custo afundado.

A base é o diferencial entre o preço físico local e a cotação futura. Para o produtor que usa o mercado futuro como proteção, o risco de preço é substituído pelo risco de base, tipicamente menor mas não nulo \cite{Hull2018}. A base do café brasileiro é negativa e sazonal: o físico negocia abaixo do futuro convertido, com diferencial mais amplo durante a colheita, quando a oferta local é abundante, e mais estreito na entressafra.

\section{Mapeamento Sistemático da Literatura}\label{sec2:msl}

Para posicionar o AgroVisible frente ao estado da arte foi conduzido um mapeamento sistemático da literatura seguindo as diretrizes de \citeonline{Petersen2015}. O mapeamento buscou responder a três questões: em que medida a literatura incorpora variáveis exógenas na previsão de preços agrícolas; quais métodos de otimização e decisão sequencial são empregados para prescrever momentos de venda pós-colheita; e como os modelos existentes tratam custos de oportunidade, armazenagem, risco de base e restrições financeiras do produtor.

A busca combinou termos de inteligência artificial (\textit{deep learning}, LSTM, \textit{neural networks}), de commodities agrícolas (\textit{coffee price}, \textit{agricultural commodities}, \textit{crop forecasting}) e de tomada de decisão (\textit{decision support}, \textit{dynamic programming}, \textit{prescriptive analytics}, \textit{optimization}) nas bases Scopus, IEEE Xplore, Periódicos CAPES e ACM Digital Library. A \autoref{tab:msl-funil} apresenta o funil de seleção.

\begin{table}[htb]
	\centering
	\caption{Funil de seleção do mapeamento sistemático.}\label{tab:msl-funil}
	\begin{tabular}{l r}
		\hline
		\textbf{Etapa} & \textbf{Registros} \\ \hline
		Scopus & 240 \\
		IEEE Xplore & 289 \\
		Periódicos CAPES & 18 \\
		ACM Digital Library & 1 \\
		\textbf{Total bruto} & \textbf{548} \\ \hline
		Excluídos após triagem de título e resumo & 353 \\
		Contexto geral de otimização agrícola (\textit{borderline}) & 15 \\
		Aceitos: modelos preditivos para preços de safras & 14 \\
		Aceitos: previsão de preço com suporte à decisão & 75 \\
		Aceitos: núcleo direto sobre café & 5 \\ \hline
		\textbf{Referências estratégicas selecionadas} & \textbf{20} \\ \hline
	\end{tabular}
	\fonte{Do autor (2026).}
\end{table}

Os critérios de exclusão removeram trabalhos sobre mineração, metais, ações, criptomoedas, petróleo e aplicações fora do agronegócio. Dos 195 registros aceitos, 20 foram selecionados como referências estratégicas por leitura completa, organizados em três pilares: núcleo direto sobre café, suporte prescritivo e otimização agrícola, e modelagem preditiva avançada com variáveis exógenas.

\subsection{Núcleo direto sobre café}

Quatro trabalhos tratam especificamente da previsão de preços de café. \citeonline{JinXu2024} aplicam processos gaussianos com otimização bayesiana a doze anos de cotações diárias, alcançando RMSE relativo de 2,05\,\%, e encerram o estudo na métrica de ajuste. \citeonline{K2023} propõem uma arquitetura híbrida de convolução com LSTM bidirecional e demonstram sua superioridade sobre modelos estatísticos em séries voláteis, mas operam sobre uma série global agregada, sem distinção entre praças físicas. \citeonline{LeNgoc2023} avaliam modelos de aprendizado de máquina no mercado vietnamita de robusta e mostram que a inclusão de variáveis exógenas, como clima e preço de combustíveis, reduz o erro de previsão de forma consistente; é o único dos quatro que explora exógenas, e o único que trata o robusta, ainda que de forma isolada. \citeonline{Aloka2026} apresentam um sistema integrado para o Sri Lanka que combina visão computacional para classificação de grãos com um módulo consultivo de mercado de caráter descritivo.

\subsection{Suporte prescritivo e otimização agrícola}

O segundo pilar reúne trabalhos que vão além da previsão. \citeonline{Mn2026} e \citeonline{Gamage2026} propõem sistemas de apoio à decisão para a comercialização pós-colheita de culturas diversas, mas a prescrição em ambos é baseada em regras sobre a previsão pontual, sem modelo formal de decisão sequencial nem contabilidade de custos de carregamento. \citeonline{Suaza2023} investigam o efeito da defasagem dos dados em um sistema de decisão para o preço do suíno, questão diretamente relevante ao alinhamento temporal adotado aqui. Os trabalhos com aprendizado por reforço, como \citeonline{Wang2024}, \citeonline{KeFan2024} e \citeonline{PerezPons2021}, aplicam Q-learning profundo a planejamento de cultivo, leilões e mercados multiagente, respectivamente; nenhum trata a decisão de venda de um estoque colhido sob custo de oportunidade. \citeonline{Bernard2025} formulam um modelo bioeconômico de otimização orientado por clima simulado para planejamento agrícola, o exemplo mais próximo de otimização formal, porém voltado à decisão de plantio e não à de comercialização.

\subsection{Modelagem preditiva com variáveis exógenas}

O terceiro pilar concentra-se em arquiteturas híbridas para previsão de preços de grãos, sobretudo milho e soja. \citeonline{Chen2025a}, \citeonline{Chen2025b}, \citeonline{Zhang2025} e \citeonline{Huang2025} combinam decomposição de sinal (VMD, EMD), seleção de atributos e redes recorrentes bidirecionais, e reportam ganhos de acurácia sobre a LSTM simples. \citeonline{Destriana2025} comparam ARIMA, LSTM e Transformers para volatilidade de preços de alimentos, concluindo que arranjos híbridos estatístico-neurais são mais robustos. \citeonline{Damanik2025} avaliam LSTM, SVR e sua combinação para o milho. Em todos, a avaliação termina em RMSE, MAE ou MAPE, e nenhum compara a previsão neural com a expectativa embutida na curva de futuros, que é a referência natural para uma commodity com mercado a termo líquido.

\section{A Lacuna Científica}\label{sec2:lacuna}

O \autoref{qua:lacuna} sintetiza a comparação entre os trabalhos do núcleo sobre café e o AgroVisible segundo cinco critérios: a cultura e a praça modeladas, o algoritmo, a modelagem de custos de carregamento, o tratamento do risco de base e a existência de prescrição formal de venda.

\begin{quadro}[htb]
	\centering
	\caption{Posicionamento do AgroVisible frente aos trabalhos do núcleo sobre café.}\label{qua:lacuna}
	\small
	\begin{tabular}{L{2.6cm} L{2.2cm} L{2.5cm} L{1.9cm} L{1.6cm} L{2.6cm}}
		\hline
		\textbf{Estudo} & \textbf{Cultura} & \textbf{Algoritmo} & \textbf{Custos} & \textbf{Base} & \textbf{Prescrição} \\ \hline
		Le Ngoc et al. (2023) & Robusta, Vietnã & LSTM, GRU, RF & Não & Não & Não \\
		K et al. (2023) & Café, global & CNN-BLSTM & Não & Não & Não \\
		Jin e Xu (2024) & Café, global & Processos gaussianos & Não & Não & Não \\
		Aloka et al. (2026) & Café, Sri Lanka & Visão computacional e ML & Não & Não & Consultiva \\
		\textbf{AgroVisible} & Arábica, Brasil (MG/SP) & LSTM e programação dinâmica & Armazenagem, quebra e Selic & CEPEA vs. ICE & Política ótima por indução retroativa \\ \hline
	\end{tabular}
	\fonte{Do autor (2026).}
\end{quadro}

Da síntese emergem quatro lacunas que este trabalho se propõe a preencher.

\paragraph{Previsão sem prescrição.} A literatura sobre café encerra-se na métrica de erro. Saber que o preço tende a subir não responde à pergunta do produtor, que é se deve vender agora, parte agora ou tudo daqui a alguns meses, e essa resposta depende da compensação entre valor esperado, dispersão e custo de carregar o estoque. Nenhum dos trabalhos mapeados formula a comercialização como problema de decisão sequencial. O AgroVisible acopla a saída da rede a um MDP resolvido por indução retroativa e entrega uma política, não uma previsão.

\paragraph{Ausência dos custos reais de carregamento.} Nenhum trabalho modela o custo de oportunidade do capital nem as tarifas de armazenagem. Uma valorização esperada de 5\,\% em três meses gera prejuízo se o custo financeiro e físico de estocar no mesmo período for de 6\,\%. O AgroVisible desconta explicitamente armazenagem, quebra técnica e Selic em cada aresta do grafo de decisão.

\paragraph{Desconsideração da base.} Os trabalhos sobre café usam séries globais ou de bolsa como se fossem o preço do produtor. O cafeicultor brasileiro recebe o preço físico local, que difere da bolsa por um diferencial sazonal e volátil. O AgroVisible modela o preço físico CEPEA/ESALQ como alvo, usa a base contra o futuro como variável de entrada e incorpora o risco de base na ação de \textit{hedge}.

\paragraph{Avaliação estatística sem validação financeira.} Todos os trabalhos mapeados avaliam a previsão por RMSE, MAE ou MAPE. Um modelo com erro médio menor pode errar a direção justamente nas reversões que mais importam para a venda física. Além disso, nenhum compara a rede com a curva de futuros, que é a expectativa que o produtor já possui sem modelo algum. O AgroVisible é avaliado, além das métricas usuais, por \textit{backtesting} de receita líquida em safras fora da amostra contra as heurísticas praticadas pelo produtor, e a rede é comparada com a curva de futuros ajustada pela base sazonal.

Uma lacuna secundária diz respeito à integração entre as duas variedades de café. \citeonline{LeNgoc2023} modelam o robusta com variáveis exógenas, mas nenhum trabalho usa o preço do robusta como preditor do arábica, apesar da substituição parcial entre elas na indústria. O AgroVisible inclui o indicador físico do robusta e a razão entre os dois preços como variáveis de entrada.

Em síntese, a contribuição deste trabalho não está em uma arquitetura neural mais acurada, mas na integração entre previsão e decisão sob uma contabilidade de custos realista, e na demonstração empírica de que essa integração pode gerar valor financeiro mesmo quando a previsão isolada não supera as referências de mercado. Os capítulos seguintes descrevem como essa hipótese foi testada.
